\documentclass[../../../include/open-logic-section]{subfiles}

\begin{document}

\olfileid{inc}{tcp}{cqn}
% Part: incompleteness
% Chapter: theories-computability
% Section: extensions-of-q-not-decidable
这一想法是通过将“可证性”概念显式化，来加强不完全性定理的证明。第一不完全性定理的证明表明，如果一个理论既可靠又完备，那么它就会证明一个实际上断言自身不可证的语句。然而，仔细检查该证明就会发现，可靠性假设可以被更弱的一致性假设所替代。（这是 Rosser（罗瑟） 在 1936 年观察到的，从而加强了 Gödel 1931 年的原始结果。）

\olsection{$\Th{Q}$ 的一致扩张是不可判定的}

\begin{explain}
回忆一下，如果一个理论对任何公式~$!A$ 都不同时证明 $!A$ 和 $\lnot !A$，那么它就是\emph{一致的}。由于从矛盾可以推出一切，不一致的理论是平凡的：每个语句都是可证的。显然，如果一个理论是 $\omega$-一致的，那么它就是一致的。但是，一致性是一个更弱的要求（即存在一致但不是 $\omega$-一致的理论）。我们可以将 \olref[oqn]{thm:oconsis-q} 中的假设减弱为简单的一致性，从而得到一个更强的定理。
\end{explain}

\begin{lem}
不存在“通用的可计算关系”。也就是说，不存在具有以下性质的二元可计算关系 $R(x,y)$：若 $S(y)$ 是一个一元可计算关系，则总存在某个 $k$，使得对每个 $y$，$S(y)$ 为真当且仅当 $R(k,y)$ 为真。
\end{lem}

\begin{proof}
假设 $R(x,y)$ 是一个通用的可计算关系。令 $S(y)$ 为关系 $\lnot R(y,y)$。
由于 $S(y)$ 可计算，对某个 $k$，$S(y)$ 等价于 $R(k,y)$。
但这样一来，$S(k)$ 就同时等价于 $R(k,k)$ 和 $\lnot R(k,k)$，矛盾。
\end{proof}


\begin{thm}
设 $\Th{T}$ 为任意包含 $\Th{Q}$ 的一致理论，则 $\Th{T}$ 是不可判定的。
\end{thm}


\begin{proof}
假设 $\Th{T}$ 是 $\Th{Q}$ 的一个一致且可判定的扩张。
我们将通过使用 $\Th{T}$ 来定义一个通用的可计算关系，从而得出矛盾。

令 $R(x,y)$ 成立当且仅当
\begin{quote}
$x$ 编码一个公式 $!D(u)$，并且 $\Th{T}$ 证明 $!D(\num y)$。
\end{quote}
由于我们假设 $\Th{T}$ 是可判定的，所以 $R$ 是可计算的。
我们来证明 $R$ 是通用的。
如果 $S(y)$ 是任意可计算关系，则它在 $\Th{Q}$ 中（因而也在 $\Th{T}$ 中）可由一个公式 $!D_S(u)$ 表示。
于是对每个 $n$，我们有
\begin{eqnarray*}
S(\num n) & \lif & T \vdash !D_S(\num n) \\
& \lif & R(\Gn{!D_S(u)}, n)
\end{eqnarray*}
以及
\begin{eqnarray*}
\lnot S(\num n) & \lif & T \vdash \lnot !D_S(\num n) \\
& \lif & T \not\vdash !D_S(\num n) \quad \text{（因为 $\Th{T}$ 是一致的）} \\
& \lif & \lnot R(\Gn{!D_S(u)}, n).
\end{eqnarray*}
也就是说，对每个 $y$，$S(y)$ 为真当且仅当
$R(\Gn{!D_S(u)}, y)$ 为真。所以 $R$ 是通用的，从而得到了我们想要的矛盾。
\end{proof}

令“真算术”为理论 $\Setabs{!A}{ \Sat{\Nat}{!A}}$，即算术语言中在标准解释下为真的所有语句的集合。

\begin{cor}
真算术是不可判定的。
\end{cor}

%In Section~\olref{undefinability:truth:section} we will state a stronger
%result, due to Tarski.

\end{document}